\section{Background, Terminology, and Research Position}
\label{sec:background}

Early malware terminology separates self-replication from host dependence and
propagation \cite{cohen1987,spafford1989}. Contemporary operational labels add
initial-access mechanism, staged-delivery role, payload objective, persistence,
and privacy impact \cite{nist80083,kharraz2016,scaife2016}. These distinctions
guide corpus design but are not interchangeable with carrier structure.
A media file may carry an exploit trigger without being a self-replicating
program. A downloader may arrive through a document, archive, executable, or a
parser flaw in nominal media. Conversely, suspicious metadata or an ambiguous
container does not establish any malware family. This study therefore reports
measurable carrier properties and keeps behavioral attribution outside the
acceptance decision.

Delivery and storage systems establish different upstream properties: sender or
server authentication, access control, transport integrity, provenance, or
confidentiality. None of these properties entails whole-file grammar,
canonicality, or safe publication. In end-to-end encrypted messaging the
inspection point must remain local because moving plaintext inspection to a
relay would alter the confidentiality model \cite{unger2015,cohngordon2020,
abelson2024}. Email, cloud synchronization, web download, and direct import
reach the same local admission problem through different upstream assumptions.
The proposed method therefore begins at the media trust boundary and treats
transport provenance as context rather than acceptance evidence.

Parser-differential research shows why one recognizer is insufficient. Different
tools can disagree about valid prefixes, terminal offsets, duplicated fields,
integer interpretation, and nested structure \cite{jana2012,momot2016}. Work on
canonical parsing, language-oriented security, and generated parsers replaces
permissive recognition with an explicit input language and complete parsing
\cite{samuel2012,nail2014,everparse2019,ipg2023}. SafeDocs applies the related
principle that complex data should be reduced to well-defined, safely processable
forms \cite{darpaSafeDocs}.
Recent studies of file ambiguity and polyglots further demonstrate that
extension, media type, and isolated parser success do not determine one
whole-file meaning \cite{koch2025,you2025}. Open reconstruction studies for
images, PDF, and OLE documents instead evaluate format-specific reduction and
fidelity \cite{belkind2023,dubin2023,dubin2024ole}.

Content disarm and reconstruction operationalizes reduction by deleting excluded
structures or by decoding and re-encoding content. The relevant security claim
depends on what is retained. Structural reconstruction parses the source and
serializes an allowed structure while retaining selected compressed payloads.
Semantic reconstruction derives a new representation from decoded pixels,
samples, frames, or equivalent media semantics. Structural reconstruction can
remove metadata, normalize layout, and exclude unaccounted suffixes. It cannot
support a claim that retained compressed payload bytes have been neutralized.
Semantic reconstruction supports a stronger source-independence claim, but only
under explicit decoder, encoder, resource, and fidelity assumptions.

Throughout, let \(\mathbb O=\{0,\ldots,255\}\) be the octet alphabet and
\(\mathbb O^*\) the set of finite byte strings. For \(b\in\mathbb O^*\),
\(|b|\) denotes octet length; \([P]\in\{0,1\}\) is the Iverson indicator of
proposition \(P\); and \(\bot\) denotes an undefined selector result or an
absent published path, according to type. Sets and finite maps are compared only
within their declared domains; all indexed sums, products, and conjunctions use
the finite index sets stated locally. One invocation fixes policy snapshot
\(p\) and parser-registry version \(v\). Subscripts retain them only where
policy comparison or evidence traceability would otherwise be ambiguous.

Let the attacker-influenced request evidence be \(r=(x,n,m)\), where \(x\) is
the input octet string, \(n\) is one portable display basename, and \(m\) is
the supplied media type. Let \(c=(p,d)\) be trusted processing context, where
\(p\) is an immutable policy snapshot and \(d\) contains the enforced deadline
and cancellation state. One processing request is therefore
\begin{equation}
  u=(r,c)=((x,n,m),(p,d)).
  \label{eq:processing-unit}
\end{equation}
An untrusted caller may request earlier cancellation or a shorter deadline but
cannot weaken \(p\) or extend the host-imposed processing budget. The request
excludes a caller-controlled output path. A final path is derived only after
acceptance from a policy-controlled directory, a sanitized basename, and the
canonical output suffix.

A format profile is more specific than a broad media family. A source profile
\(s\) fixes the admissible input evidence and reconstruction route. An output
profile \(k\) fixes a canonical output grammar, media type, suffix set, resource
limits, and parser version. Policy also registers the permitted
source-to-output mappings and the minimum reconstruction level. This distinction
is necessary because strict routes may convert WAV, FLAC, or Ogg input to MP3,
and WebM input to MP4. Static and animated WebP are separate profiles because
their component languages and reconstruction claims differ. Likewise, WebM and
general Matroska cannot share one native structural route merely because both
use EBML.

For policy $p$ and profile $k$, let $L_{p,k}\subseteq\mathbb O^{*}$ denote the
canonical output language. A sound authorizing parser starts at offset zero,
terminates at $|y|$, admits every component at its parsed location, and checks
ordering, cardinality, cross-reference, integrity, and resource constraints.
If \(\rho_{p,v,k}(y)\) is its report under policy \(p\) and parser-registry
version \(v\), the soundness direction used here is
\begin{equation}
  \mathsf{CompleteOK}_{p}(\rho_{p,v,k}(y),k,y)
  \Longrightarrow y\in L_{p,k}.
  \label{eq:complete-profile}
\end{equation}
The reverse implication is not assumed: a conservative parser may reject a
language member. A successful prefix parse is not complete acceptance, and
membership is not permanent across policy or parser-registry versions.
Canonicality always refers to a named profile and frozen policy.

Reconstruction levels form the ordered set
\begin{equation}
  \NotProcessed \prec \Structural \prec \Semantic.
  \label{eq:reconstruction-order}
\end{equation}
The order describes admissible claim strength, not a universal quality ranking.
A policy requirement \(q_{p,s}\) for source profile \(s\) is satisfied only
when the performed level obeys \(\ell\geq q_{p,s}\). Unsupported input is blocked
rather than misrepresented as structurally reconstructed.

Verdicts are separated from pipeline errors:
\begin{equation}
\begin{aligned}
  V&=\{\Clean,\Suspicious,\Blocked\},\\
  O&=\mathsf{Result}(V\times \mathcal E,\mathsf{PipelineError}),
\end{aligned}
  \label{eq:verdict-space}
\end{equation}
where $\mathcal E$ is the set of bounded stage-evidence records. In enforcement mode, a clean verdict
means that every mandatory check succeeded and no policy warning or blocker
remains. A suspicious verdict means that no blocking condition remains but at
least one policy warning does. A blocked verdict records at least one blocking
condition. Suspicious and blocked results are both non-deliverable. A pipeline
error is also non-deliverable but remains distinct for diagnosis and
availability analysis. These cases are mutually exclusive, and no failure is
converted into a clean result by exception handling or fallback routing.

An authorizing parser is considered separate from construction when it is
invoked after reconstruction through a code path distinct from the handler's
source parser and writer. This separation does not imply that the parser was
independently verified. The stronger term \emph{implementation-diverse
validator} is reserved for external parser or decoder families used in the
differential evaluation. The distinction prevents separation within one
implementation from being reported as independent replication.

\begin{table*}[t]
\centering
\caption{Core notation used by the formal method and evidence model.}
\label{tab:core-notation}
\small
\begin{tabularx}{\textwidth}{L{0.16\textwidth}L{0.24\textwidth}Y}
\toprule
Scope & Symbol & Meaning \\
\midrule
Request & \(u=((x,n,m),(p,d))\) & Bytes, portable name, supplied type, immutable policy, deadline, and cancellation context \\
Registry & \(s,k;\ R_p,\mathsf{Route}_p\) & Source/output profiles, eligible set, and fail-closed selector \\
Policy & \(\ell,q_{p,s};\ \Omega_p\) & Performed/required level and source--level--output relation \\
Policy & \(L_{p,k},\Gamma_{p,k}\) & Canonical byte language and location-aware component policy \\
Candidate & \(a=(y,n_o,m_o,k_o,\ell,e)\) & Reconstructed bytes, generated evidence, performed level, and bounded stage record \\
Validation & \(\rho_j,\Pi^{\mathrm{auth}}_{p,v,k}\) & Parser report and configured authorizing-parser set \\
Decision & \(\Accept_{p,v},\mathsf{Pub}_{p,v}\) & Candidate authorization and non-overwriting publication transition \\
Study & \(\mathbb G,G_{\mathrm{sci}}\) & Three-valued gate domain and aggregate scientific gate \\
\bottomrule
\end{tabularx}
\end{table*}

\begin{table*}[t]
\centering
\caption{Direct overlap between representative prior-work lines and obligations
F1--F8. ``Evidence'' means that a line can test or motivate an obligation but
does not establish it as a publication-boundary guarantee.}
\label{tab:research-position}
\small
\begin{tabularx}{\textwidth}{L{0.27\textwidth}L{0.20\textwidth}Y}
\toprule
Representative line & Direct F1--F8 overlap & Remaining composition gap \\
\midrule
Signature and malware scanning \cite{nist80083} & Supporting evidence for F8 & No direct F1--F7 guarantee; a clean scan does not authorize grammar or reconstruction \\
Canonical/generated parsing \cite{samuel2012,nail2014,everparse2019,ipg2023} & F1--F2 & Does not jointly enforce F3--F8 at reconstruction and publication boundaries \\
Content disarm and reconstruction \cite{darpaSafeDocs,belkind2023,dubin2023,dubin2024ole} & Format-dependent F2/F4 evidence & Transformation success alone does not establish F1, F3, F5--F8 \\
Parser differentials and polyglot analysis \cite{jana2012,koch2025,you2025} & Evidence for F1/F3 & Detects disagreement but does not supply reconstruction, bounds, or clean-only publication \\
This work & F1--F8 as one authorization relation & Guarantees remain conditional on registered-parser, construction, and filesystem assumptions \\
\bottomrule
\end{tabularx}
\end{table*}

\Cref{tab:research-position} does not claim that the cited systems are weaker
within their own scopes. It identifies the combination studied here: unique
routing, explicit reconstruction strength, complete validation of emitted
bytes, resource control, and publication restricted to clean verdicts, with a
distinct falsifier for every obligation.
